% =============================================================================
% CAPÍTULO 5: PROCESSOS TÉCNICOS
% Alinhado à ISO/IEC/IEEE 12207:2021 (6.4) e MPS.BR
% Fonte: manual-processos/manual/capitulos/05-processos-tecnicos.md
% =============================================================================

\aberturacapitulo{Capítulo 5}{Processos Técnicos}

\secao{Processo de definição de requisitos}

A Engenharia de Requisitos da \termochave{CineSoft Solutions} segue o processo formal da \termochave{ISO/IEC/IEEE 12207:2021 (6.4.2 e 6.4.3)}, garantindo que as necessidades dos clientes e regulamentos legais (como a LGPD) sejam traduzidas em requisitos verificáveis.

\subsecao{Catálogo Formal de Requisitos Funcionais e Não Funcionais}

\tabelaeditorial{Catálogo de Requisitos Centrais do Sistema CineDoom}{Engenharia CineSoft (2026)}{tab:requisitos_sistema}{p{2.2cm} p{4.2cm} X p{2.2cm}}{
  \rowcolor{BrandTableHeader}
  \textbf{\textcolor{BrandNavy}{Código}} & \textbf{\textcolor{BrandNavy}{Nome do Requisito}} & \textbf{\textcolor{BrandNavy}{Descrição e Critério de Aceitação}} & \textbf{\textcolor{BrandNavy}{Classificação}} \\
  \midrule
  \codigoinline{RF-001} & Catálogo de Filmes e Sessões & Gerenciamento de filmes em cartaz e horários sem choque de grade na mesma sala. & Funcional \\
  \codigoinline{RF-002} & Trava Atômica de Assento & Bloqueio exclusivo temporário por 10 min (\codigoinline{SeatHold}) com prevenção de double-booking. & Funcional \\
  \codigoinline{RF-003} & Baixa Presencial na Bilheteria & Registro e conferência presencial de pagamento no guichê físico com transação atômica. & Funcional \\
  \codigoinline{RF-004} & Despacho de Vouchers por E-mail & Envio automatizado de voucher com instruções e chave de reserva via Resend API. & Funcional \\
  \codigoinline{RNF-001} & Concorrência Transacional & Garantia estrita de consistência ACID em cenários de alta concorrência de pré-vendas. & Não Funcional \\
  \codigoinline{RNF-002} & Privacidade e Proteção LGPD & Anonimização de dados pessoais, consentimento explícito e revogação de cadastro. & Não Funcional \\
  \codigoinline{RNF-003} & Tempo de Resposta de Checkout & Tempo total de validação e bloqueio de assento inferior a 400ms sob carga normal. & Não Funcional \\
}

\secao{Processo de arquitetura e design do sistema}

A arquitetura do CineDoom adota o padrão de \destaqueazul{Monólito Modular com Clean Architecture}, desacoplando regras de negócio corporativas de detalhes de infraestrutura e bibliotecas externas.

% FIGURA COM PRINT DA BILHETERIA IMPORTADA DE ASSETS
\figuracomfonte{prints/print-bilheteria-atendente.png}{0.88\textwidth}{Interface Operacional do Atendente: Guichê Físico e Mapa de Assentos}{CineSoft Solutions (2026)}{fig:bilheteria_atendente}

\subsecao{Diagrama de Contêineres (C4 Model Nível 2)}

O sistema é composto por:
\begin{listapontos}
  \itemponto{Single Page / Server Component App}{Interface Next.js 16 servida na borda pela Vercel, proporcionando renderização híbrida rápida para mobile e guichê.}
  \itemponto{API Transacional de Negócios}{Camada de serviços escrita em TypeScript que orquestra a máquina de estados de reserva e venda.}
  \itemponto{Repositório Relacional Serverless}{Banco de dados Neon PostgreSQL conectado via Prisma ORM com suporte a transações ACID concorrentes.}
  \itemponto{Provedor de E-mail Externo}{Serviço Resend integrado via HTTPS REST para despacho assíncrono de notificações transacionais.}
\end{listapontos}

% FIGURA COM PRINT DO PORTAL DO CLIENTE IMPORTADA DE ASSETS
\figuracomfonte{prints/print-portal-cliente.png}{0.88\textwidth}{Portal do Cliente Web: Consulta de Grade e Reserva de Poltronas}{CineSoft Solutions (2026)}{fig:portal_cliente}

\secao{Processo de implementação e integração de software}

A implementação do CineDoom orienta-se pela manutenibilidade, cobertura de testes e aderência aos padrões de código:
\begin{listapontos}
  \itemponto{Padrões de Código}{Padronização rigorosa com ESLint e Prettier aplicada automaticamente via Git Hooks (Husky).}
  \itemponto{Migrations de Banco}{Todas as alterações de schema são versionadas no Prisma e aplicadas via pipeline automatizada no deploy de homologação.}
  \itemponto{Tratamento de Exceções}{Mapeamento de erros de concorrência com respostas padronizadas em JSON para o cliente web.}
\end{listapontos}

\secao{Processos de verificação, validação e transição (V\&V)}

Em conformidade com a ISO 12207 (6.4.9 e 6.4.10), a CineSoft estabelece distinção clara entre Verificação e Validação:

\tabelaeditorial{Matriz de Verificação vs. Validação do CineDoom}{CineSoft Solutions (2026)}{tab:verif_valid}{p{4.2cm} p{5.5cm} X}{
  \rowcolor{BrandTableHeader}
  \textbf{\textcolor{BrandNavy}{Dimensão}} & \textbf{\textcolor{BrandNavy}{Verificação (O produto foi construído corretamente?)}} & \textbf{\textcolor{BrandNavy}{Validação (O produto certo foi construído?)}} \\
  \midrule
  Objetivo Principal & Avaliar se os artefatos de código cumprem rigorosamente as especificações técnicas e normas. & Avaliar se a solução atende às necessidades reais dos operadores e do público nos cinemas. \\
  Técnicas Adotadas & Testes unitários com Vitest, análise estática de tipos no TypeScript e code reviews formais. & Testes de aceitação em ambiente de homologação (\textit{staging}), testes de usabilidade e piloto com cinemas. \\
  Critério de Aprovação & 100\% dos testes automatizados verdes e zero vulnerabilidades críticas detectadas no pipeline. & Assinatura formal do Termo de Aceite pelos gerentes de cinema e aprovação dos fluxos operacionais. \\
}

\secao{Processos de operação, manutenção e descontinuação}

A operação do CineDoom inclui monitoramento de saúde do banco de dados, gestão de backups diários e procedimentos de rollback imediato em caso de incidentes.

\verificacaoconformidade{Ponto de Controle — Processos Técnicos}{%
  \begin{enumerate}[label=\textbf{\alph*)}, leftmargin=0.8cm, itemsep=3pt]
    \item Como a arquitetura modular do CineDoom garante que falhas no envio de e-mails pelo Resend não interrompam a venda física de ingressos na bilheteria?
    \item Quais são os critérios de entrada e saída para que uma nova versão do sistema avance do ambiente de homologação para produção?
  \end{enumerate}
}

\notatecnica{%
  \itemcaixa{O isolamento multi-tenant do CineDoom utiliza cláusula obrigatória de \codigoinline{tenantId} em todas as consultas Prisma, garantindo separação lógica absoluta entre os dados de diferentes redes exibidoras.}
}

\saibamaisnormativo{%
  \itemsaibamais{Consulte as cláusulas 6.4.4 (\textit{Architecture Definition}), 6.4.9 (\textit{Verification}) e 6.4.10 (\textit{Validation}) da norma \linkcaixa{https://www.iso.org/standard/63711.html}{ISO/IEC/IEEE 12207:2021} para detalhamento dos processos de ciclo de vida técnico.}
}

\sintesecapitulo{%
  \itemcaixa{A engenharia do CineDoom comprova que arquitetura de alta concorrência e processos de ciclo de vida rigorosos são indispensáveis para garantir a integridade de vendas de ingressos digitais.}
}

